\documentclass[10pt]{scrartcl}
\usepackage[mdthm,colorsec,noauthor,noasy]{evan}
\usepackage{handout}
\newcommand{\aug}[2]{\left[\begin{array}{@{}#1@{}}#2\end{array}\right]}
\newcommand{\vv}[1]{\mathbf{#1}}
\DeclareMathOperator{\spn}{Span}
\title{Math 54 --- Week 2}
\subtitle{Solutions}
\author{}\date{}\ihead{}
\begin{document}
\maketitle

\begin{prob}\textbf{P1.}\par\idea{Row reduce $[\,\vv a_1\ \vv a_2\ \vv a_3\ \vv b\,]$; the solution gives the weights.}
\[ \aug{rrr|r}{1&2&1&5\\0&1&-1&1\\2&3&4&12}\xrightarrow{R_3-2R_1}\aug{rrr|r}{1&2&1&5\\0&1&-1&1\\0&-1&2&2}\xrightarrow{R_3+R_2}\aug{rrr|r}{1&2&1&5\\0&1&-1&1\\0&0&1&3} \]
$x_3=3$, $x_2=1+3=4$, $x_1=5-8-3=-6$. So $\vv b=-6\vv a_1+4\vv a_2+3\vv a_3$. Check: $(-6,0,-12)+(8,4,12)+(3,-3,12)=(5,1,12)$.
\end{prob}

\begin{prob}\textbf{P2.}\par\idea{$\vv y$ is in the span exactly when $[\,\vv v_1\ \vv v_2\ \vv y\,]$ is consistent.}
\[ \aug{rr|r}{1&-3&h\\0&1&-5\\-2&8&-3}\xrightarrow{R_3+2R_1}\aug{rr|r}{1&-3&h\\0&1&-5\\0&2&2h-3}\xrightarrow{R_3-2R_2}\aug{rr|r}{1&-3&h\\0&1&-5\\0&0&2h+7} \]
Consistent exactly when $2h+7=0$: $h=-7/2$.
\end{prob}

\begin{prob}\textbf{P3.}\par\idea{Columns span $\RR^3$ exactly when there is a pivot in every row.}
\begin{steps}
  \ii $A$: $R_2+4R_1$ gives $(0,14,10)$; $R_3+3R_1$ gives $(0,7,5)$; then $R_3-\tfrac12R_2$ gives a zero row. Only two pivots: the columns do \emph{not} span $\RR^3$.
  \ii $B$: $R_3-2R_1$ gives $(0,0,1)$. Echelon form $\begin{bmatrix}1&2&0\\0&1&3\\0&0&1\end{bmatrix}$: three pivots, one in every row. The columns span $\RR^3$.
\end{steps}
\end{prob}

\begin{prob}\textbf{P4.}
Definition: $5\begin{bmatrix}2\\0\\1\end{bmatrix}-2\begin{bmatrix}-1\\3\\4\end{bmatrix}=\begin{bmatrix}10+2\\0-6\\5-8\end{bmatrix}=\begin{bmatrix}12\\-6\\-3\end{bmatrix}$. Row-vector rule: $2\cdot5+(-1)(-2)=12$; $0\cdot5+3(-2)=-6$; $1\cdot5+4(-2)=-3$. Same answer.
\end{prob}

\begin{prob}\textbf{P5.}\par\idea{Reduce, solve for basic variables, split by free variable.}
\[ \begin{bmatrix}1&3&-3&7\\0&1&-4&5\\2&7&-10&19\end{bmatrix}\to\begin{bmatrix}1&0&9&-8\\0&1&-4&5\\0&0&0&0\end{bmatrix} \]
(Steps: $R_3-2R_1$ gives $(0,1,-4,5)$, then $R_3-R_2$ gives a zero row; finally $R_1-3R_2$.) $x_1=-9x_3+8x_4$, $x_2=4x_3-5x_4$, $x_3,x_4$ free:
\[ \vv x=x_3\begin{bmatrix}-9\\4\\1\\0\end{bmatrix}+x_4\begin{bmatrix}8\\-5\\0\\1\end{bmatrix}. \]
\end{prob}

\begin{prob}\textbf{P6.}\par\idea{Particular solution plus the homogeneous solutions.}
\[ \aug{rrr|r}{1&2&-1&1\\2&5&1&4\\3&7&0&5}\to\aug{rrr|r}{1&2&-1&1\\0&1&3&2\\0&1&3&2}\to\aug{rrr|r}{1&0&-7&-3\\0&1&3&2\\0&0&0&0} \]
$x_1=-3+7x_3$, $x_2=2-3x_3$:
\[ \vv x=\begin{bmatrix}-3\\2\\0\end{bmatrix}+x_3\begin{bmatrix}7\\-3\\1\end{bmatrix}. \]
A line through $(-3,2,0)$ parallel to $(7,-3,1)$; the solution set of the homogeneous system is the parallel line through $\vv0$. Check: $A(-3,2,0)=(1,4,5)$ and $A(7,-3,1)=(0,0,0)$.
\end{prob}

\begin{prob}\textbf{P7.}
\begin{steps}
  \ii (a) False. A $3\times2$ matrix has at most 2 pivots, so not one in every row.
  \ii (b) True. By the theorem on solution sets, the solutions are $\vv p+\vv v_h$ with $\vv v_h$ a solution of $A\vv x=\vv0$; the only such $\vv v_h$ is $\vv0$.
  \ii (c) True. At most 3 pivots among 5 columns, so at least 2 free variables; the homogeneous system is consistent, so it has nontrivial solutions.
\end{steps}
\end{prob}

\begin{prob}\textbf{P8.}\par\idea{One equation, one basic variable $x_1$, two free variables.}
$x_1=4-2x_2+3x_3$:
\[ \vv x=\begin{bmatrix}4\\0\\0\end{bmatrix}+x_2\begin{bmatrix}-2\\1\\0\end{bmatrix}+x_3\begin{bmatrix}3\\0\\1\end{bmatrix}. \]
The homogeneous equation has solution set $\spn\{(-2,1,0),(3,0,1)\}$, a plane through $\vv0$. The given equation's solution set is that plane shifted by $(4,0,0)$: a parallel plane not through $\vv0$.
\end{prob}

\end{document}
